\textbf{Phần 1. Trắc nghiệm lựa chọn}
\setcounterpageref{ex}{0}

\begin{ex} %[1D10V39]
Tính: $35 + 42$.
\choice
{\True $77$}
{$76$}
{$78$}
{$75$}
\loigiai{
Ta có: $35 + 42 = 77$.
\textbf{Đáp số:} $77$.\par
\textbf{Đáp án của câu hỏi là A}
}
\end{ex}

\begin{ex} %[1D10V39]
Tính: $89 - 54$.
\choice
{$34$}
{$36$}
{\True $35$}
{$33$}
\loigiai{
Ta có: $89 - 54 = 35$.
\textbf{Đáp số:} $35$.\par
\textbf{Đáp án của câu hỏi là C}
}
\end{ex}

\begin{ex} %[1D10V39]
Tính: $40 + 50 - 30$.
\choice
{$50$}
{\True $60$}
{$70$}
{$59$}
\loigiai{
Thực hiện: $40 + 50 = 90$; $90 - 30 = 60$.
\textbf{Đáp số:} $60$.\par
\textbf{Đáp án của câu hỏi là B}
}
\end{ex}

\begin{ex} %[1D10V39]
Điền số thích hợp: $24 + \dots = 67$.
\choice
{$42$}
{$44$}
{$41$}
{\True $43$}
\loigiai{
Số cần điền là: $67 - 24 = 43$.
\textbf{Đáp số:} $43$.\par
\textbf{Đáp án của câu hỏi là D}
}
\end{ex}

\begin{ex} %[1D10V39]
Điền số thích hợp: $98 - \dots = 55$.
\choice
{\True $43$}
{$42$}
{$44$}
{$41$}
\loigiai{
Số cần điền là: $98 - 55 = 43$.
\textbf{Đáp số:} $43$.\par
\textbf{Đáp án của câu hỏi là A}
}
\end{ex}

\begin{ex} %[1D10V39]
Tìm $x$ biết: $x - 32 = 45$.
\choice
{$76$}
{$78$}
{\True $77$}
{$75$}
\loigiai{
Ta có: $x = 45 + 32 = 77$.
\textbf{Đáp số:} $77$.\par
\textbf{Đáp án của câu hỏi là C}
}
\end{ex}

\begin{ex} %[1D10V39]
Tìm $y$ biết: $y + 14 = 88$.
\choice
{$73$}
{\True $74$}
{$75$}
{$72$}
\loigiai{
Ta có: $y = 88 - 14 = 74$.
\textbf{Đáp số:} $74$.\par
\textbf{Đáp án của câu hỏi là B}
}
\end{ex}

\begin{ex} %[1D10V39]
Lớp 1A có 34 học sinh, lớp 1B có 35 học sinh. Hỏi cả hai lớp có bao nhiêu học sinh?
\choice
{$68$}
{$70$}
{$67$}
{\True $69$}
\loigiai{
Tổng số học sinh cả hai lớp là: $34 + 35 = 69$ (học sinh).
\textbf{Đáp số:} $69$.\par
\textbf{Đáp án của câu hỏi là D}
}
\end{ex}

\begin{ex} %[1D10V39]
Một trang trại có 86 con thỏ, đã bán được 43 con thỏ. Hỏi trang trại còn lại bao nhiêu con thỏ?
\choice
{\True $43$}
{$42$}
{$44$}
{$41$}
\loigiai{
Số thỏ còn lại là: $86 - 43 = 43$ (con thỏ).
\textbf{Đáp số:} $43$.\par
\textbf{Đáp án của câu hỏi là A}
}
\end{ex}

\begin{ex} %[1D10V39]
Một cửa hàng buổi sáng bán được 25 chiếc áo, buổi chiều bán được 31 chiếc áo. Cả ngày bán được bao nhiêu chiếc áo?
\choice
{$55$}
{$57$}
{\True $56$}
{$54$}
\loigiai{
Cả ngày bán được: $25 + 31 = 56$ (chiếc áo).
\textbf{Đáp số:} $56$.\par
\textbf{Đáp án của câu hỏi là C}
}
\end{ex}

\begin{ex} %[1D10V39]
Tìm số tròn chục lớn nhất có hai chữ số trừ đi 40.
\choice
{$40$}
{\True $50$}
{$60$}
{$49$}
\loigiai{
Số tròn chục lớn nhất là 90: $90 - 40 = 50$.
\textbf{Đáp số:} $50$.\par
\textbf{Đáp án của câu hỏi là B}
}
\end{ex}

\begin{ex} %[1D10V39]
Tìm số bé nhất có hai chữ số cộng với 65.
\choice
{$74$}
{$76$}
{$73$}
{\True $75$}
\loigiai{
Số bé nhất có hai chữ số là 10: $10 + 65 = 75$.
\textbf{Đáp số:} $75$.\par
\textbf{Đáp án của câu hỏi là D}
}
\end{ex}

\textbf{Phần 2. Trắc nghiệm đúng sai}
\setcounterpageref{ex}{0}

\begin{ex} %[1D10V39]
Xét tính đúng hoặc sai của các khẳng định toán học sau đây:
\choiceTF
{\True Kết quả của bài toán «Hiệu giữa số lớn nhất có hai chữ số và số 55 bằng bao nhiêu» là $44$.}
{\True Kết quả của bài toán «Bác An thu hoạch được 54 bắp ngô, bác Ba thu hoạch được 45 bắp ngô. Cả hai bác thu hoạch được bao nhiêu bắp ngô» là $99$.}
{Kết quả của bài toán «Bác An thu hoạch nhiều hơn bác Ba bao nhiêu bắp ngô» là $8$.}
{Giá trị của phép tính $12 + 23 + 34 + 20$ là $88$.}
\loigiai{
\begin{itemize}
    \item Ý a) Đúng: Số lớn nhất có hai chữ số là 99: $99 - 55 = 44$.
\textbf{Đáp số:} $44$.
    \item Ý b) Đúng: Cả hai bác thu hoạch được: $54 + 45 = 99$ (bắp ngô).
\textbf{Đáp số:} $99$.
    \item Ý c) Sai: Kết quả đúng là $9$, không phải $8$. Số bắp ngô bác An nhiều hơn bác Ba là: $54 - 45 = 9$ (bắp ngô).
\textbf{Đáp số:} $9$.
    \item Ý d) Sai: Kết quả đúng là $89$, không phải $88$. Thực hiện: $12 + 23 = 35$; $35 + 34 = 69$; $69 + 20 = 89$.
\textbf{Đáp số:} $89$.
\end{itemize}
\begin{itemize}
    \item a) \textbf{Đúng}.
    \item b) \textbf{Đúng}.
    \item c) \textbf{Sai}.
    \item d) \textbf{Sai}.
\end{itemize}
}
\end{ex}

\begin{ex} %[1D10V39]
Xét tính đúng hoặc sai của các bài toán và phép tính sau đây:
\choiceTF
{Giá trị của phép tính $98 - 30 - 25$ là $42$.}
{Kết quả của bài toán «Một cây vải dài 95 mét, cửa hàng đã bán đi 63 mét. Cây vải còn lại bao nhiêu mét» là $33$.}
{\True Kết quả của bài toán «Một đoạn dây điện dài 78 cm. Bác thợ cắt bớt 25 cm. Đoạn dây điện còn lại dài bao nhiêu xăng-ti-mét» là $53\text{ cm}$.}
{\True Kết quả của bài toán «Bác An thu hoạch nhiều hơn bác Ba bao nhiêu bắp ngô» là $9$.}
\loigiai{
\begin{itemize}
    \item Ý a) Sai: Kết quả đúng là $43$, không phải $42$. Thực hiện: $98 - 30 = 68$; $68 - 25 = 43$.
\textbf{Đáp số:} $43$.
    \item Ý b) Sai: Kết quả đúng là $32$, không phải $33$. Số mét vải còn lại là: $95 - 63 = 32$ (mét).
\textbf{Đáp số:} $32$.
    \item Ý c) Đúng: Đoạn dây còn lại dài: $78 - 25 = 53$ (cm).
\textbf{Đáp số:} $53$.
    \item Ý d) Đúng: Số bắp ngô bác An nhiều hơn bác Ba là: $54 - 45 = 9$ (bắp ngô).
\textbf{Đáp số:} $9$.
\end{itemize}
\begin{itemize}
    \item a) \textbf{Sai}.
    \item b) \textbf{Sai}.
    \item c) \textbf{Đúng}.
    \item d) \textbf{Đúng}.
\end{itemize}
}
\end{ex}

\textbf{Phần 3. Trả lời ngắn}
\setcounterpageref{ex}{0}

\begin{ex} %[1D10V39]
Tìm số tự nhiên $m$ biết rằng: $m + 21 = 55 + 12$.
\shortans{46}
\loigiai{
\textbf{Đáp số:} $46$\par
Vế phải: $55 + 12 = 67$. Ta có: $m + 21 = 67 \Rightarrow m = 67 - 21 = 46$.
\textbf{Đáp số:} $46$.
}
\end{ex}

\begin{ex} %[1D10V39]
Có bao nhiêu số tròn chục lớn hơn 30 và bé hơn 90?
\shortans{5}
\loigiai{
\textbf{Đáp số:} $5$\par
Các số tròn chục đó là: $40, 50, 60, 70, 80$. Có tất cả $5$ số.
\textbf{Đáp số:} $5$.
}
\end{ex}

\begin{ex} %[1D10V39]
Hai số có tổng bằng 88, số bé là 35. Hỏi số lớn là số mấy?
\shortans{53}
\loigiai{
\textbf{Đáp số:} $53$\par
Số lớn là: $88 - 35 = 53$.
\textbf{Đáp số:} $53$.
}
\end{ex}

\begin{ex} %[1D10V39]
Hai số có hiệu bằng 24, số bé là 45. Hỏi số lớn là số mấy?
\shortans{69}
\loigiai{
\textbf{Đáp số:} $69$\par
Số lớn là: $45 + 24 = 69$.
\textbf{Đáp số:} $69$.
}
\end{ex}

\begin{ex} %[1D10V39]
Tìm số có hai chữ số giống nhau biết số đó lớn hơn 70 và bé hơn 80.
\shortans{77}
\loigiai{
\textbf{Đáp số:} $77$\par
Số đó là số $77$.
\textbf{Đáp số:} $77$.
}
\end{ex}

\begin{ex} %[1D10V39]
Tính tổng các số tròn chục từ 10 đến 40.
\shortans{100}
\loigiai{
\textbf{Đáp số:} $100$\par
Tổng là: $10 + 20 + 30 + 40 = 100$.
\textbf{Đáp số:} $100$.
}
\end{ex}
